% Standalone problem document, generated from the adjacent README.md.
% Compile with XeLaTeX. Mathematical content is maintained in Markdown.
\documentclass[11pt,a4paper]{article}
\usepackage[fontsize=10.5pt]{scrextend}
\usepackage[margin=22mm,headheight=15pt,headsep=9mm,footskip=11mm]{geometry}
\usepackage{fontspec}
\setmainfont{texgyrepagella}[Extension=.otf,UprightFont=*-regular,BoldFont=*-bold,ItalicFont=*-italic,BoldItalicFont=*-bolditalic]
\setsansfont{texgyreheros}[Extension=.otf,UprightFont=*-regular,BoldFont=*-bold,ItalicFont=*-italic,BoldItalicFont=*-bolditalic]
\setmonofont{texgyrecursor}[Extension=.otf,UprightFont=*-regular,BoldFont=*-bold,ItalicFont=*-italic,BoldItalicFont=*-bolditalic,Scale=MatchLowercase]
\usepackage{amsmath,amssymb}
\usepackage{unicode-math}
\setmathfont{texgyrepagella-math.otf}
\usepackage{xcolor,microtype,enumitem,needspace,fancyhdr}
\usepackage{bookmark}
\definecolor{ink}{HTML}{17344B}
\definecolor{accent}{HTML}{197D80}
\definecolor{muted}{HTML}{596773}
\hypersetup{colorlinks=true,linkcolor=accent,urlcolor=accent,pdftitle={MI-32 - Spectral
norms of independent entries with regular moment
growth},pdfauthor={Open Problems in Numerical Linear Algebra}}
\urlstyle{same}
\setlength{\parindent}{0pt}
\setlength{\parskip}{5pt plus 1pt minus 1pt}
\linespread{1.04}
\setlength{\emergencystretch}{3em}
\widowpenalty=10000
\clubpenalty=10000
\displaywidowpenalty=10000
\allowdisplaybreaks[1]
\setlist{nosep,leftmargin=1.5em}
\providecommand{\tightlist}{\setlength{\itemsep}{0pt}\setlength{\parskip}{0pt}}
\providecommand{\pandocbounded}[1]{#1}
\pagestyle{fancy}
\fancyhf{}
\fancyhead[L]{\sffamily\footnotesize\color{muted}OPEN PROBLEMS IN NUMERICAL LINEAR ALGEBRA}
\fancyhead[R]{\sffamily\bfseries\footnotesize\color{accent}MI-32}
\fancyfoot[L]{\parbox[b]{0.9\textwidth}{\sffamily\fontsize{8}{10}\selectfont\color{muted}Editorial ratings; a bounded search cannot certify that a problem remains unsolved.\\Literature check: 2026-09-11}}
\fancyfoot[R]{\sffamily\footnotesize\color{muted}\thepage}
\renewcommand{\headrulewidth}{0.3pt}
\renewcommand{\headrule}{\hbox to\headwidth{\color{accent}\leaders\hrule height \headrulewidth\hfill}}
\makeatletter
\renewcommand{\section}{\Needspace{5\baselineskip}\@startsection{section}{1}{0pt}{12pt plus 2pt minus 2pt}{4pt}{\sffamily\bfseries\large\color{ink}}}
\renewcommand{\subsection}{\Needspace{4\baselineskip}\@startsection{subsection}{2}{0pt}{9pt plus 1pt minus 1pt}{3pt}{\sffamily\bfseries\normalsize\color{ink}}}
\makeatother
\setcounter{secnumdepth}{0}
\begin{document}
{\sffamily\small\color{muted}Matrix inequalities and norms\par}
\vspace{3pt}
{\raggedright\hyphenpenalty=10000\exhyphenpenalty=10000\sffamily\bfseries\fontsize{21}{24}\selectfont\color{ink}Spectral
norms of independent entries with regular moment growth\par}
\vspace{7pt}
{\sffamily\small\textbf{Difficulty:} challenging\quad\textbf{Importance:} interesting
to the community\par}
{\sffamily\small\bfseries\color{accent}Status: Partially resolved\par}
\vspace{4pt}
{\color{accent}\hrule height 0.8pt}
\vspace{6pt}
\textbf{Topic:} Expected spectral norms of inhomogeneous random matrices

\textbf{Rating rationale:} The missing upper estimate must account for
localized random fluctuations across all variance profiles. Even
weighted random signs retain an iterated-logarithm gap. A sharp estimate
would improve quantitative understanding of the largest singular value
in random matrix models.

\subsection{Statement}\label{statement}

For a real random variable \(Z\) and \(r>0\), write
\(\|Z\|_{L_r}=(\mathbb E|Z|^r)^{1/r}\). Fix \(\alpha\ge1\). Let
\(X=(X_{ij})\in\mathbb R^{n\times n}\) have independent mean-zero
entries with finite absolute moments of every order, satisfying

\[\|X_{ij}\|_{L_{2r}}\le\alpha\|X_{ij}\|_{L_r}
\qquad(1\le i,j\le n,\ r\ge1).\]

Set \([n]=\{1,\ldots,n\}\) and

\[
M(X)=\max_i\left(\sum_j\mathbb E X_{ij}^2\right)^{1/2}
+\max_j\left(\sum_i\mathbb E X_{ij}^2\right)^{1/2},
\]

\[
D(X)=\max_{1\le k\le n}\;
\min_{\substack{I\subseteq[n]\\|I|\le k}}
\sup_{\substack{s,t\in\mathbb R^n\\\|s\|_2,\|t\|_2\le1}}
\left\|\sum_{\substack{i\notin I\\j\notin I}}X_{ij}s_it_j
\right\|_{L_{\ln(k+1)}}.
\]

\textbf{Conjecture (Latała--Świątkowski).} For every \(\alpha\ge1\)
there is a constant \(C_\alpha>0\) such that, for every \(n\ge1\) and
every such matrix,

\[\mathbb E\|X\|_2\le C_\alpha\bigl(M(X)+D(X)\bigr),\]

where \(\|X\|_2\) is the spectral norm. The constant is independent of
the dimension and entry laws. The same deterministic index set \(I\)
deletes both rows and columns; the supremum is outside the
random-variable moment. The definition includes \(L_{\ln2}\) when
\(k=1\), using the displayed moment functional even though its exponent
is below one.

The reverse inequality up to a constant depending only on \(\alpha\) is
proved, so the target is equivalent to the source's two-sided
comparison. No identical-distribution or symmetry assumption is imposed
on the entries.

\newpage
\subsection{Known cases and numerical
significance}\label{known-cases-and-numerical-significance}

The source proves the Gaussian-mixture case under the same moment
condition. Weighted signs \(X_{ij}=a_{ij}\varepsilon_{ij}\), with
independent fair \(\varepsilon_{ij}\in\{-1,1\}\), satisfy the condition
with \(\alpha=1\) and form the source's earlier Conjecture 1.2. Latała
proves this case when \(a_{ij}\in\{0,1\}\) and proves the general
weighted-sign estimate with an additional factor of order
\(\log\log\log n\). Meller's 2026 result gives an iterated-logarithm
loss for a wider class of symmetric entries; it does not remove that
loss. These cases remain grouped in this entry.

The spectral norm is the largest singular value and measures the largest
Euclidean amplification by a random matrix. The formula combines row and
column variance scales with moments of bilinear forms after deleting a
limited number of coordinates, accounting for concentration on small
parts of a matrix that simpler variance-only bounds can miss.

\subsection{References and status
check}\label{references-and-status-check}

\begin{itemize}
\tightlist
\item
  R. Latała and W. Świątkowski, \emph{Norms of randomized circulant
  matrices}, Electronic Journal of Probability 27 (2022), paper 80,
  1--23. \href{https://doi.org/10.1214/22-EJP799}{DOI};
  \href{https://arxiv.org/pdf/2106.03139v2}{current arXiv v2}, dated
  2022-05-27. Conjecture 4.3 on preprint p.25, with condition (25) on
  p.24, states the target. Theorem 4.1 supplies the lower bound;
  Proposition 4.4 and the paragraph before it cover Gaussian mixtures.
  Conjecture 1.2 on p.2 is the weighted-sign special case.
\item
  R. Latała, \emph{On the spectral norm of Rademacher matrices}.
  \href{https://doi.org/10.1090/tran/9637}{DOI};
  \href{https://arxiv.org/html/2405.13656v2}{current arXiv v2}, dated
  2025-08-18. Equations (1.2)--(1.3), Theorem 1.1, and Theorem 1.9 give
  the weighted-sign conjecture and the stated partial results. This
  paper uses the comparable truncated logarithm \(\max\{1,\ln k\}\) in
  place of \(\ln(k+1)\).
\item
  R. Meller, \emph{Spectral norm of matrices with independent entries up
  to polyloglog},
  \href{https://arxiv.org/html/2512.23673v2}{arXiv:2512.23673v2}, dated
  2026-01-29, introduction, equation (3), and Theorem 1.1. It explicitly
  identifies the remaining logarithmic gap and Conjecture 4.3.
\end{itemize}

On 2026-09-11, checked the original paper's current version, both later
papers, and exact-title, author/conjecture, Rademacher spectral-norm,
proof/counterexample and 2026 searches. No full resolution was found.
Gaussian operator-norm results and the general Bernoulli-process theorem
do not prove this displayed formula. This is a bounded literature check.
\end{document}
